\chapter{Предисловие}

\begin{center} 
	Здравствуй, мой дорогой читатель!
\end{center}

Пособие, которое ты сейчас читаешь, "--- результат титанического труда многих людей. И мы все верим, что результат получился более чем достойный. 

Цель данной книги – облегчить подготовку студентов Московского Физико"--~Технического Института к устному выпускному квалификационному государственному экзамену по математике, попросту к ГОСу. Как вы уже поняли, оглавление представляет собой программу к ГОСу 2016 года. Заметьте, что оно <<кликабельно>> в pdf-версии данной книги, что упрощает работу с книгой. И <<кликабельно>> не только оглавление данной книги, но и всевозможные числа и названия, указывающие на теоремы, которые уже использовались ранее в книге, а также объекты в мини-содержаниях перед главами. Надеюсь, это кому-нибудь поможет. 

Надо бы отметить для любителей учебников и лекций определённых авторов, что мы писали билеты, существенно опираясь на учебные пособия, лекции различных преподавателей. Список соответствия билетов из программы и названий курсов от кафедры высшей математики материалам, которыми мы в основном пользовались (подчеркиваю, в основном) смотрите ранее, в списке литературы.

Мы были предельно внимательны к составлению данной книги, стараясь уменьшить количество опечаток и повысить качество излагаемого материала. Но мы отказываемся от ответственности за всевозможные недочеты в этой книге (пожалуй, главный недочет этой книги "--- чрезмерная избыточность в некоторых местах, а иногда, наоборот, недостаточно материала), ведь мы, на данный момент, всего лишь студенты, а главное "--- люди, которые могут ошибаться. И поэтому, прошу вас не забывать отправлять нам (ссылки ниже) сообщения о любых неточностях, опечатках, ошибках, недочетах. Также пишите, если хотите дать совет или выразить любые личные пожелания. Вместе с вами можно довести эту книгу до очень хорошего пособия.

Не могу не отметить доброжелательного и внимательного отношения всех студентов МФТИ к этому пособию. Хочу сказать всем, кто присылал сообщения об опечатках и ошибках: ``Спасибо''. Также хочется выразить особую благодарность Кудашову Аркадию, Лузянину Артемию, Проскину Роману, Вербе Глебу и Браславскому Илье за непосредственное соучастие в написании этой книги и выразить признательность Брицыну Евгению и Дроботу Олегу за многочисленные комментарии и исправления.

Не обошлось даже без участия преподавательского состава МФТИ. Так, например, Максим Широбоков, преподаватель теории вероятностей, случайных процессов и математической статистики, прочитал билеты по теории вероятностей, оставил важные и ценные комментарии, и в ходе долгой дискуссии после редактирования билетов одобрил последние. За это все я очень благодарен ему.

Также Чубаров Игорь Андреевич на своей очной консультации упомянул, что читал это пособие, и сказал, что билеты по его предмету (аналитическая геометрия и линейная алгебра) написаны хорошо.
 
Мне лишь остается выразить надежду, что настоящее пособие поможет студентам при изучении математики в целом и подготовке к ГОСу. Но все же я настоятельно рекомендую пользоваться не только данным пособием при подготовке к ГОСу.
\hyphenation{ГОСе}

\vspace*{0.7\baselineskip} 

\textit{Желаю всем отличных результатов~на~ГОСе.}

\mbox{}

\noindent\textit{Диденко А.А.}

\vspace*{0.1\baselineskip} 
\noindent\href{https://vk.com/didenko_andre}{$
	\begin{array}{l}
	\includegraphics[width=\iconwidth]{pictures/reference/vk}
	\end{array}$/didenko\_andre}

\noindent\href{https://telegram.me/didenko_andre}{$
	\begin{array}{l}
	\includegraphics[width=\iconwidth]{pictures/reference/telegram}
	\end{array}$/didenko\_andre}
    
\noindent\href{https://instagram.com/didenko_andre}{$
    \begin{array}{l}
    \includegraphics[width=\iconwidth]{pictures/reference/instagram}
    \end{array}$/didenko\_andre}

\vspace*{0.5\baselineskip} 

\noindent PS. Пользуйтесь Дидодичкой на здоровье, с удовольствием и умом. :)

\vspace*{0.5\baselineskip} 

\noindent PPS. Вы можете сами нам помочь, сделав \textit{Pull Request} или указав на существующий \textit{Issue} в следующем репозитории:

\vspace*{0.1\baselineskip}
 
\noindent\href{https://github.com/DidenkoAndre/GOSBook}{$
	\begin{array}{l}
	\includegraphics[width=\iconwidth]{pictures/reference/git}
	\end{array}$/didenkoandre/GOSBook}

\vspace*{0.5\baselineskip} 

\noindent PPPS. Подпишитесь на наш Telegram-бот с помощью команды \textsf{/subscribe}, и вы будете получать свежую версию автоматически!

\vspace*{0.1\baselineskip} 

\noindent\href{https://t.me/GOSBook_Bot}{$
	\begin{array}{l}
	\includegraphics[width=\iconwidth]{pictures/reference/telegram}
	\end{array}$/GOSBook\_Bot}

\vspace*{0.5\baselineskip} 

\iffalse
\noindent $\bullet$ Напоследок, вы можете поддержать проект материально! Перевести средства на развитие проекта можно на счет {\large 5244 6873 8510 3816} \textsf{(MasterCard)}.
\fi

\vspace*{0.5\baselineskip} 

\textit{Спасибо!}
